\section*{الفصل \ref{ch:b3:stem-cells} --- الخلايا الجذعية، والتجدد، والشيخوخة}

\begin{solution}{exo:b3:stem-cells:1}
\omterm{def:b3:stem-cells:stem}{التجدد الذاتي}: انقسام ينتج بنتًا واحدة على الأقل لا تزال
\omterm{def:b3:stem-cells:stem}{خلية جذعية} --- أي خلايا \emph{Lgr5} في قاعدة الخبيئة.
والقدرة: مدى المصائر التي لا تزال الخلية تقدر على اتخاذها ---
\omterm{def:b3:stem-cells:stem}{فالخلية الجذعية} في الخبيئة متعددة القدرات، تصنع خلايا
امتصاصية، وكأسية، ومعوية صماوية، وخلايا بانيث، وخلايا
خصلية. \omterm{def:b3:stem-cells:stem}{والحاضنة}: خلايا بانيث وإشاراتها Wnt وNotch وEGF، التي
تثبّت الجذعية في مكانها. \omterm{def:b3:stem-cells:stem}{وخلية العبور المضخِّمة}: بنت ملتزمة
تنقسم أربع مرات أو خمسًا أخرى في الخبيئة قبل أن تتمايز في
طريقها إلى أعلى الزغابة.
\end{solution}

\begin{solution}{exo:b3:stem-cells:2}
أنتجت خلايا نقي مفردة مستعمرات طحالية تحوي كريات حمرًا
ومحببات وصفيحات: أي خلية واحدة، وسلالات عدة --- فتعدد
القدرات، وبيّنت الاستجابة الخطية للجرعة أن كل مستعمرة جاءت
من خلية واحدة. وأما \omterm{def:b3:stem-cells:stem}{التجدد الذاتي} فاقتضى فأرًا ثانيًا:
فخلايا من مستعمرة، محقونةً في متلقٍّ مشعع آخر، صنعت مستعمرات
جديدة، فتكون الخلية المؤسِّسة قد أنتجت بنات بقدرتها هي.
\end{solution}

\begin{solution}{exo:b3:stem-cells:3}
\emph{Oct4}، و\emph{Sox2}، و\emph{Klf4}، و\emph{c-Myc}. \omterm{def:b3:chromatin-epigenetics:nucleosome}{وصبغين}
الأرومة الليفية قد أغلق مورّثات \omterm{def:b3:stem-cells:pluripotency}{كمال القدرة} خلف \omterm{def:b3:chromatin-epigenetics:nucleosome}{صبغين متغاير}
\omterm{def:b3:chromatin-epigenetics:methylation}{ومثيلة DNA}، فتجد العوامل مواقع متاحة قليلة في البداية وعليها
أن تجنّد إنزيمات إعادة تشكيل على مدى دورات خلوية كثيرة؛ وعلى
الخلية أيضًا أن تطفئ برنامجها هي، وأن تفلت من الشيخوخة، وأن
تمر بحالة وسيطة عشوائية يتعثر فيها معظم الخلايا أو يموت.
فلا تخرج كاملةَ القدرة إلا الخلية النادرة التي تنجح فيها كل
خطوة، بعد أسابيع.
\end{solution}

\begin{solution}{exo:b3:stem-cells:4}
تحتفظ البلاناريا بخلايا جذعية بالغة كاملة القدرات، هي
النيوبلاستات، التي تهاجر إلى الجرح وتبني كل ما ينقص. وأما
السمندر فلا مجمع كهذا لديه: بل تزيل خلايا متمايزة عند
الجذمور تمايزها إلى سلائف، تكوّن مع خلايا جذعية مقيمة في
النسيج \omterm{def:b3:stem-cells:regeneration}{بلاستيما} تعيد تشغيل برنامج الطرف. وقد أظهرت طعوم
كراغل لنسج موسومة مفردة أن خلايا \omterm{def:b3:stem-cells:regeneration}{البلاستيما} تبقى وفيّة
لأصلها --- فالعضل يصنع عضلًا، والغضروف غضروفًا --- \omterm{def:b3:stem-cells:regeneration}{فالبلاستيما}
خليط من سلائف محصورة السلالة، لا كتلة كاملة القدرة.
\end{solution}

\begin{solution}{exo:b3:stem-cells:5}
$(11 - 5)/0.075 = 80$ مضاعفة. ومن \qty{8}{kb}: $(8 - 5)/0.075 = 40$.
ومع تعويض \qty{60}{\%} يكون الفقد الصافي \qty{30}{bp}: أي $6000/30
= 200$.
\end{solution}

\begin{solution}{exo:b3:stem-cells:6}
أربعة عشر انقسامًا جذعيًا في اليوم تعطي 28 بنتًا؛ ويجب أن
تبقى 14 خلايا جذعية، فتلتزم 14 --- أي بنت ملتزمة واحدة لكل
\omterm{def:b3:stem-cells:stem}{خلية جذعية} في اليوم. وكل ملتزمة تصير $2^{4} = 16$ خلية: أي
$14\times 16 = 224$ خلية زغابية في اليوم، وهي من رتبة
\num{300} المرصودة (وانقسام عبور خامس يعطي \num{448}). وتسهم
طبقة الخلايا الجذعية بأربعة عشر انقسامًا، وتسهم طبقة العبور
بالخلايا \num{200} الباقية.
\end{solution}

\begin{solution}{exo:b3:stem-cells:7}
$2\times 10^{11}/2^{20} = 1.9\times 10^{5}$ بنت ملتزمة في اليوم؛
ومع \num{10000} \omterm{def:b3:stem-cells:stem}{خلية جذعية}، 19 لكل \omterm{def:b3:stem-cells:stem}{خلية جذعية} في اليوم --- في
مقابل انقسام واحد في السنة. والعشرون انقسامًا محسوبة لكل
سلالة من \omterm{def:b3:stem-cells:stem}{الخلية الجذعية} إلى الكرية الحمراء، غير أن السلائف
الوسيطة ليست لمرة واحدة: فمجامع السلائف المتعددة القدرات
والملتزمة تُديم نفسها أسابيع بانقساماتها هي، فيقع الانقسام
كله تقريبًا في طبقات السلائف ولا تُستدعى \omterm{def:b3:stem-cells:stem}{الخلية الجذعية} إلا
نادرًا.
\end{solution}

\begin{solution}{exo:b3:stem-cells:8}
$\mu(50) = 10^{-3}\mathrm{e}^{1.74} = 5.7\times 10^{-3}$؛ و$\mu(70) =
10^{-3}\mathrm{e}^{3.48} = 0.032$؛ و$\mu(90) = 10^{-3}\mathrm{e}^{5.22} =
0.19$ في السنة. ومع $\mu_{0}/\gamma = 0.0115$: $S(70) = \exp[-0.0115
\times 31.5] = 0.70$، و$S(90) = \exp[-0.0115\times 184] = 0.12$. والوسيط:
$30 + \ln(61)/0.087 = 77$. وبتنصيف $\mu_{0}$: $\ln(121)/0.087 = 55$،
فالوسيط 85 --- أي زمن مضاعفة واحد، ثماني سنوات، مكسوب. وبتنصيف $\gamma$
إلى $0.0435$: $\ln(31)/0.0435 = 79$، فالوسيط 109: أي إن إبطاء الأس
يكسب أربعة أضعاف ما يكسبه تنصيف الثابت.
\end{solution}

\begin{solution}{exo:b3:stem-cells:9}
على الثلث الباقي أن يتثلث: أي $\log_{2}3 = 1.6$ انقسام لكل خلية.
وهو \omterm{def:b3:stem-cells:regeneration}{نمو تعويضي}، لأن الفصوص الموجودة تتضخم بانقسام الخلايا في
موضعها؛ وأما الفصوص المزالة، بقنواتها وأوعيتها ومعمارها،
فلا يُعاد بناؤها --- فالكبد يستعيد كتلته ووظيفته لا شكله.
\end{solution}

\begin{solution}{exo:b3:stem-cells:10}
في كل حدث ينتقل حجم النسيلة $k$ إلى $k + 1$ أو $k - 1$ باحتمال
متساوٍ، فلا يتغير حجمها المتوقع أبدًا؛ وحين ينتهي المشي يكون
$N$ (استيلاء) أو $0$ (ضياع)، والتوقع $N\,P + 0 = k$
يعطي $P = k/N$. وخلية واحدة: $1/14$. والطفرة المحايدة في \omterm{def:b3:stem-cells:stem}{خلية
جذعية} واحدة تتثبت في خبيئتها باحتمال $1/14$ وتضيع في غضون
أشهر فيما عدا ذلك --- فالانجراف ينقّي ثلاث عشرة طفرة من كل
أربع عشرة.
\end{solution}

\begin{solution}{exo:b3:stem-cells:11}
صار المشي منحازًا: فالنسيلة تكسب مكانًا أكثر مما تخسر
مكانًا، فينمو حجمها المتوقع ويرتفع احتمال الاستيلاء نحو
اليقين مع نمو $k$ --- فلأفضلية \qty{10}{\%} في
\omterm{def:b3:stem-cells:stem}{حاضنة} من أربع عشرة، نحو ثلاثة أضعاف $1/14$ لخلية واحدة،
وأعلى من ذلك متى حملت النسيلة بضعة أماكن. وعلى مدى عقود
تستولي الخلايا الجذعية ذات هذه الطفرات (\emph{DNMT3A}،
و\emph{TET2} في النقي) على حاضناتها، وعند السبعين يكون خُمس
الناس حصة كبيرة من دمهم من نسيلة واحدة. وليس ذلك سرطانًا:
فالخلايا لا تزال تتمايز وتطيع تراتبها. غير أن الطفرة التالية
صار لها نسيلة من مليارات لتنشأ فيها، وارتفع خطر الابيضاض
عشرة أضعاف.
\end{solution}

\begin{solution}{exo:b3:stem-cells:12}
الوسيط $= 30 + \gamma^{-1}\ln(1 + \gamma\ln 2/\mu_{0})$: فعند $\mu_{0} =
10^{-2}$، $30 + \ln 7/0.087 = 52$؛ وعند $10^{-3}$: 77؛ وعند $10^{-4}$: $30 +
\ln 604/0.087 = 104$. وفي صورة غومبرتز الصرفة يضيف كل خفض عشري
$\ln 10/\gamma = 26$ سنة نفسها. وتتناقص العوائد عمليًا
لأن المكاسب التاريخية جاءت من إزالة حد مستقل عن العمر
--- الخمج، والولادة، والحوادث، أي ثابت ميكهام $A$ في $\mu = A
+ \mu_{0}\mathrm{e}^{\gamma t}$ --- ومتى صار $A$ صغيرًا في مقابل
الأسّي عند الأعمار التي تهم، لم يكد خفضه أكثر يغيّر شيئًا؛
وأما $\mu_{0}$ الجوهري فقد تحرك أقل بكثير. وخفض \qty{20}{\%}
في $\gamma$ (إلى $0.070$) يعطي $30 + \ln 49/0.070 = 86$: أي تسع سنوات،
مكسوبةً في كل عمر لا بتأجيل بضعة أسباب للموت.
\end{solution}

\begin{solution}{pb:b3:stem-cells:1}
\textbf{1.} $3.5\times 10^{11}$ في اليوم، و$4\times 10^{6}$ في الثانية.
\textbf{2.} $2^{20} \approx 10^{6}$ خلية لكل بنت ملتزمة؛
و$3.5\times 10^{11}/10^{6} = 3.3\times 10^{5}$ بنت ملتزمة في اليوم.
\textbf{3.} $33$ لكل \omterm{def:b3:stem-cells:stem}{خلية جذعية} في اليوم، و\num{12000} في السنة.
\textbf{4.} على مجامع السلائف أن تُديم نفسها: فالسلائف
المتعددة القدرات والملتزمة تنقسم أسابيع وتحفظ أعدادها هي،
فيقع كل انقسام تقريبًا من العشرين في طبقات \omterm{def:b3:stem-cells:regeneration}{تجدد} نفسها، ولا
تسهم \omterm{def:b3:stem-cells:stem}{الخلية الجذعية} ببنت ملتزمة إلا نادرًا --- من رتبة مرة
في السنة.
\textbf{5.} $10 - 20\times 0.06 = \qty{8.8}{kb}$: أي أعلى بكثير من
$L_{\text{crit}}$، والكرية الحمراء لا تنقسم أبدًا بعد ذلك؛
فالاحتياطي لا يهم إلا السلائف التي يُعاد استعمالها.
\textbf{6.} الفقد الصافي $0.3\times 60 = \qty{18}{bp}$؛ و$6000/18 = 333$
انقسامًا؛ وعند واحد في السنة، $333$ سنة --- أي أبعد بكثير من
عمر، فالتيلوميرات لا تحدّ \omterm{def:b3:stem-cells:stem}{الخلية الجذعية} نفسها؛ وإنما يحدّها
أذى آخر.
\textbf{7.} $14$ انقسامًا في اليوم؛ و$7$ تستبدل خلايا جذعية
ضائعة و$7$ تلتزم؛ فيكون $7\times 2^{4} = 112$ خلية زغابية في اليوم.
\textbf{8.} $N^{2} = 196$ حدثًا لكل خلية عند واحد في اليوم: أي نحو
\qty{200}{d}، ستة أشهر إلى سبعة --- وهي المدة التي صارت فيها
خبيئات الكونفيتّي أحادية اللون.
\textbf{9.} $1/14 \approx \qty{7}{\%}$؛ و$10^{7}/14 \approx 7\times 10^{5}$
خبيئة.
\textbf{10.} الانجراف يطرح ثلاث عشرة طفرة من كل أربع عشرة في
غضون أشهر. وأما \omterm{def:b3:stem-cells:stem}{الانقسام اللامتناظر} مدى الحياة فكان سيبقي كل
\omterm{def:b3:stem-cells:stem}{خلية جذعية} طافرة مدى العمر، فتتراكم الطفرات في كل سلالة بلا
أن تُنقّى أبدًا.
\textbf{11.} الجذعية موضع في \omterm{def:b3:stem-cells:stem}{الحاضنة} ومجموعة إشارات، لا خاصية
دائمة لخلية: فأي خلية تبلغ إشارات بانيث تقدر على أخذ الدور.
\textbf{12.} استحثّ وسمًا متعدد الألوان في الخلايا الجذعية في
لحظة واحدة وتابع الخبيئات. فالانجراف المتناظر: تصير الخبيئات
أحادية اللون في غضون أشهر. \omterm{def:b3:stem-cells:stem}{والانقسام اللامتناظر}: تحفظ كل
خبيئة ألوانها كلها إلى ما لا نهاية، كل لون شريطًا رفيعًا
صاعدًا في الزغابة.
\textbf{13.} $6000/60 = 100$ مضاعفة.
\textbf{14.} نحو $75$ تبقى: فالخلية تعدّ الانقسامات لا الزمن
التقويمي، إذ لم يكلفها العقد في المجمّدة شيئًا.
\textbf{15.} $(8000 - 4000)/30 = 133$ سنة بعد العشرين --- أي عند
العمر 153. فمتوسط طول تيلوميرات الكريات البيض لا يحدّ العمر
بنفسه؛ وإنما تهم أقصر التيلوميرات وسائر المعالم أكثر.
\textbf{16.} \omterm{def:b3:cancer-biology:telomerase}{التيلوميراز} يزيل حاجزًا واحدًا، هو \omterm{thm:b3:stem-cells:hayflick}{الشيخوخة
التضاعفية}، ولذلك يوجد في كل \omterm{def:b3:cancer-biology:hallmarks}{سرطان} تقريبًا، إذ يحتاج إلى
انقسام بلا حد؛ غير أن الأرومات الليفية المعطاة تيلوميرازًا
حفظت نقاط تفتيشها، وتثبيطها بالتماس، وp53 سليمًا، فلم يجعلها
الانقسام غير المحدود وحده سرطانات. فهو لازم لا كافٍ.
\textbf{17.} $\mu(40) = 2.4\times 10^{-3}$، و$\mu(60) = 0.014$، و$\mu(80) =
0.077$، و$\mu(100) = 0.44$ في السنة. و$S = \exp[-0.0115(\mathrm{e}^{\gamma
(t-30)} - 1)]$: أي $0.98$ عند 40، و$0.87$ عند 60، و$0.42$ عند 80، و$0.006$ عند
100.
\textbf{18.} الوسيط $30 + \ln 61/0.087 = 77$. ويبقى عشرة في المئة عند
$\mathrm{e}^{\gamma(t-30)} = 1 + \gamma\ln 10/\mu_{0} = 201$: أي $t = 30 +
5.3/0.087 = 91$.
\textbf{19.} \qty{0.67}{mm/d}؛ و$\log_{2}1000 = 10$ مضاعفات في 60
يومًا، أي واحدة كل ستة أيام.
\textbf{20.} نحو $10^{8}$ انقسام، أي $0.1$ طفرة مسرطِنة في كل
\omterm{def:b3:stem-cells:regeneration}{تجدد} --- أي واحدة في كل عشرة أطراف. والأكسولوتلات لا تكاد
تصاب بالأورام: فهي باردة الدم بطيئة، ذات كبت ورمي متين،
وخلايا بلاستيماها تبقى محصورة السلالة وتعيد التمايز بدل أن
تواصل الانقسام؛ بل إن النسيج المتجدد يكبت الأورام المستحثة.
وأما الثديي فله خلايا أكثر، وحرارة ومعدل استقلاب أعلى، وعمر
أطول تقدر فيه نسيلة مفلتة على النمو.
\textbf{21.} انزع عصب الجذمور وازرع حبيبة تطلق البروتين
المرشَّح، أو اطعم نسيجًا مكيَّفًا بالعصب: فاستعادة \omterm{def:b3:stem-cells:regeneration}{التجدد}
تعني عاملًا قابلًا للانتشار. والمرشحون المعيَّنون في السمادر
والأكسولوتل: بروتين التدرج الأمامي nAG، والنيوريغولين-1،
وعوامل FGF.
\textbf{22.} \omterm{prop:b3:stem-cells:crypt}{إشارة Wnt} مرتفعة عند الذنب ومنخفضة عند الرأس،
والجرح عند الطرف الأمامي للقطعة يصنع رأسًا عادةً لأن Wnt
هناك منخفض؛ وخفضُ مثبط Wnt يرفع Wnt في كل مكان، فيقرأ
الجرحان ``خلفي''، فينمو ذنبان. والقطعة تعرف محورها هي من
التدرج المتبقي في عضلها، والجرح يعيد إقامة التدرج بالنسبة
إليه.
\textbf{23.} على الثلث الباقي أن يتثلث: فتنقسم كل خلية كبدية
مرة (فتبلغ الثلثين) وينقسم نصفها مرة أخرى --- أي متوسط $1.6$ انقسام.
فتتضخم الفصوص الناجية إلى الكتلة الأصلية؛ وأما الفصوص
المفقودة وقنواتها وأوعيتها فلا يُعاد بناؤها.
\textbf{24.} \omterm{def:b3:stem-cells:regeneration}{التجدد} الذي يستبدل النسيج البالي والمتأذي كان
سيبطئ التراكم الذي يقيسه الأس، فيخفض $\gamma$ ويفلطح المنحنى
--- والأكسولوتلات تُظهر شيخوخة ضئيلة. وأما الثمن فحدّ ثابت
أعلى من السرطانات في جسد حار كبير طويل العمر.
\textbf{25.} عدّة هايفليك نحو $100$ مضاعفة؛ \omterm{def:b3:stem-cells:stem}{وللخلية الجذعية}
المعانة \omterm{def:b3:cancer-biology:telomerase}{بالتيلوميراز} نحو $330$ انقسام، أي أكثر من ثلاثة قرون
عند واحد في السنة؛ والعمر الوسيط $77$ سنة.
\end{solution}
